\section{Две скорости и мост $\speedC=\caKappaGeom \caSignalSpeed$}
\label{sec:two-c}

На гексагональном срезе $(2{+}1)$ коэффициент
\[
\caKappaGeom_{\mathrm{hex}}=\sqrt{3}/2
\]
(\cref{sec:hex-slice}); на \latGCK
\[
\caKappaGeom_{\mathrm{FCC}}=1/\sqrt{2}
\]
(\cref{sec:fcc-kappa}; индекс $\mathrm{FCC}$ --- тип Браве из \cref{def:packing-types}).
Величина $\caKappaGeom$ --- отношение радиусов вписанной и описанной сфер однотактового тела.
Тождество $\speedC=\caKappaGeom \caSignalSpeed$ служит проверкой согласованности макроописания
после фиксации $\htick=\caKappaGeom t_P$ (\cref{fig:carrier-two-speeds-micro,fig:carrier-two-speeds-macro}).

\begin{definition}[Две скорости и continuum-упаковка $t_P$]
\label{def:two-speeds-formal}
\begin{itemize}
  \item \textbf{Скорость звена на $\Mlayer$:}
  $\caSignalSpeed:=\caStepSpace/\htick$ (\cref{def:minkowski-carrier-split}).
  \item \textbf{Макроскопическая скорость на $\Tlayer$:}
  $\speedC:=\caKappaGeom\,\caSignalSpeed$ с якорем CODATA после SI-моста (\cref{def:si-bridge}).
  \item \textbf{Continuum-упаковка:}
  определение $t_P:=\ell_P/\speedC$ в рамке одной скорости $\speedC$ на всех масштабах
  и подстановка $t_P$ вместо $\htick$ при чтении $\caSignalSpeed=\ell_P/t_P$.
\end{itemize}
\end{definition}

\begin{proposition}[Каузальный предел одного такта на $\Mlayer$]
\label{prop:m-one-tick-cone}
За один такт $t\to t+1$ закон~$g$ не переносит информацию за пределы каузальной звезды $N(x)$
(аксиома~$A_1$, \cref{ax:A1}): для дискретного фронта
$\|\Delta x\|_\infty\le \caStepSpace$.
Следовательно, максимальная скорость сигнала на $\Mlayer$ за такт равна $\caSignalSpeed$;
превышение за один $\htick$ без нарушения $A_1$ невозможно.
\end{proposition}

\begin{proposition}[Отношение $\caSignalSpeed/\speedC$ при FCC]
\label{prop:c0-over-c}
При $\caKappaGeom_{\mathrm{FCC}}=1/\sqrt{2}$ и \cref{def:si-bridge}:
\[
  \htick=\caKappaGeom\, t_P,\qquad
  \speedC=\caKappaGeom\,\caSignalSpeed,\qquad
  \frac{\caSignalSpeed}{\speedC}=\frac{1}{\caKappaGeom}=\sqrt{2}.
\]
Величина $\caKappaGeom$ задаётся геометрией однотактового тела (\cref{sec:fcc-kappa}),
не отношением $\speedC/\caSignalSpeed$ как свободным параметром.
\end{proposition}

\begin{proposition}[Предел скорости наблюдаемого фронта на $\Tlayer$]
\label{prop:t-speed-limit}
Скорость распространения EM-вакуумного фронта, извлекаемая прибором на $\Tlayer$
после readout $\mathcal{B}$ на масштабах $\gg\caStepSpace$,
согласована с $\speedC$ (\cref{sec:c-light-cone,ch:macro}; сверка --- CE-M-01).
Утверждение $\caSignalSpeed>\speedC$ в SI \emph{не} определяет на $\Tlayer$
состояние с групповой скоростью $v>\speedC$:
$\caSignalSpeed$ относится к паре $(\caStepSpace,\htick)$ на $\Mlayer$,
$\speedC$ --- к изотропному макро-якорю $\caKappaGeom\,\caSignalSpeed$.
\end{proposition}

\begin{corollary}[Нет кинематического доступа к $\Mlayer$]
\label{cor:no-m-kinematic-access}
По \cref{cor:no-continuum,sec:no-continuum-carrier} шаг $\caStepSpace=\ell_P$ не является
пределом $h\to 0$ для масштаба тела или прибора на $\Tlayer$.
Отображение «уменьшение линейного размера $\Rightarrow$ замена $\speedC$ на $\caSignalSpeed$''
не входит в класс допустимых морфизмов модели: наблюдатель остаётся на $\Tlayer$
с конусом $\speedC$; кинематика $\caSignalSpeed$ --- свойство обновления на $\carrierLattice\times\mathbb{Z}$.
\end{corollary}

\begin{remark}[Сверка]
\label{rem:two-c-burden}
\cref{prop:c0-over-c} --- CE-M-00 (\texttt{Kappa\_bottom\_up});
\cref{prop:t-speed-limit} --- CE-M-01 (suite \texttt{t\_macro}).
Ложное тождество «$\speedC$ на шаге $\ell_P$ при $t_P=\ell_P/\speedC$'' есть частный случай
неверной continuum-упаковки из \cref{def:two-speeds-formal}.
\end{remark}

\begin{figure}[htbp]
  \centering
  \begin{tikzpicture}[carrier panel, x=2.1cm, y=2.1cm]
  \def\tlim{1.25}
  \def\xlim{1.35}
  \draw[carrier muted, dashed] (0,0) -- (\xlim,\tlim);
  \draw[carrier muted, dashed] (0,0) -- (-\xlim,\tlim);
  \draw[carrier object] (0,0) -- (\xlim,0) node[midway, below, inner sep=1pt] {$\caStepSpace$};
  \draw[carrier object] (0,0) -- (0,\tlim) node[midway, left, inner sep=1pt] {$\caStepTime$};
  \foreach \i in {0,1,2,3} {
    \filldraw[carrier object] ({0.33*\i}, {0.33*\i}) circle (1.6pt);
  }
\end{tikzpicture}

  \caption{Микро: дискретный шаг по решётке, тактовая скорость $\caSignalSpeed$.}
  \label{fig:carrier-two-speeds-micro}
\end{figure}

\begin{figure}[htbp]
  \centering
  \begin{tikzpicture}[carrier panel, x=2.1cm, y=2.1cm]
  \def\tlim{1.25}
  \def\xlim{1.35}
  \def\caSignalSpeed{1.08}
  \def\kappa{0.7071}
  \pgfmathsetmacro{\cmacro}{\caSignalSpeed*\kappa}
  \draw[carrier muted, dashed] (0,0) -- (\xlim,\tlim);
  \draw[carrier muted, dashed] (0,0) -- (-\xlim,\tlim);
  \draw[carrier muted, dashed, line width=0.7pt] (0,0) -- (\xlim,\xlim/\caSignalSpeed);
  \draw[carrier object] (0,0) -- (\xlim,\xlim/\cmacro);
  \node[anchor=west, carrier muted, inner sep=1pt] at (0.55,0.62) {$\caSignalSpeed$};
  \node[anchor=west, inner sep=1pt] at (0.55,0.88) {$c=\symKappa \caSignalSpeed$};
  \node[anchor=north west, carrier muted, font=\scriptsize, inner sep=1pt] at (0.05,0.95)
    {$\symKappa=1/\sqrt{2}$};
\end{tikzpicture}

  \caption{Макро: осреднённый фронт, $\speedC=\caKappaGeom \caSignalSpeed$.}
  \label{fig:carrier-two-speeds-macro}
\end{figure}

